%LTex: language=pl
W ramach tej pracy przeprowadzono szereg powiązanych ze sobą eksperymentów, pozwalających na analizę właściwości modeli ASR w ramach wykorzystania ich jako część potoku przetwarzania w systemach wyszukiwania informacji za pomocą sygnału mowy. Rozdział 3 koncentrował się na właściwościach przestrzeni latentnych obecnych w enkoderach dostępnych modeli automatycznego rozpoznawania mowy w świetle enkodowania semantyki danych wejściowych, jak i możliwości zastosowania dodatkowego przetwarzania danych enkodowanych przez te modele. Wnioski pozyskane z tych eksperymentów posłużyły w doborze odpowiednich enkoderów audio w ramach treningu proponowanych w tej pracy modeli, \textit{Attention Pool Model} i \textit{Embedding Alignment Model}, które wytrenowane zostały na zbiorze AdmedVOICE w celu semantycznego wyszukiwania dokumentów tekstowych za pomocą sygnału mowy. Rozdział 5 zawierał testy modeli w ramach różnych scenariuszy testowych, oferując porównanie różnych podejść, wykorzystujących zarówno proponowane architektury, jak i klasyczne podejście kaskadowe oparte na transkrypcji sygnału mowy.

\section{Pytania badawcze}
\subsection{Czy reprezentacje ukryte generowane przez modele ASR zawierają wystarczającą informację semantyczną do realizacji zadań wyszukiwania informacji?}
W wyżej wymienionej formie, eksperymenty w tej pracy wykazały częściową zdolność wykorzystania modeli ASR w celu wyszukiwania informacji podobnych semantycznie w ramach domeny ogólnej. Wyniki badań w Rozdziale 3, wskazują na dosyć silną zdolność dyskryminacji między nagraniami mowy podobnymi semantycznie, a różnymi, w szczególności dla modeli SeamlessM4T i Whisper. Pośrednim potwierdzeniem tego wniosku jest stabilność treningu modeli \textit{Attention Pool Model} i \textit{Embedding Alignment Model}, wykorzystujących enkodery wyżej wymienionych modeli ASR, choć nie oznacza to bezpośrednio nacechowania semantycznego danych przetworzonych przez takie enkodery (Rozdział 4). Jednakże, na jednoznaczne potwierdzenie tej tezy nie pozwalają wyniki uzyskane w ramach eksperymentów zawartych w Rozdziale 5, gdzie proponowane w tej pracy architektury nie były w stanie dorównać wyszukiwaniu kaskadowemu lub modelowi embeddingowemu SONAR. Nie może oznaczać to jednak braku przydatności architektur \textit{APM} i \textit{EAM} w ramach wyszukiwania, ponieważ najprawdopodobniej najważniejszym czynnikiem prowadzącym do wyraźnie gorszych rezultatów wyszukiwania dla tych modeli jest ilość danych wykorzystanych w ramach treningu, będąca w niektórych przypadkach $10000$ razy mniejsza niż dla gotowych modeli ewaluowanych w ramach porównania \cite{m4t}. Podsumowując, nie udało się jednoznacznie udowodnić zdatności proponowanych architektur modeli wykorzystujących enkodery modeli ASR w ramach wyszukiwania informacji, lecz w pracy wykazano wiele przesłanek pozwalających sądzić, że przy większej ilości danych treningowych, architektury te mogłyby dobrze radzić sobie w ramach tego zadania.
\subsection{Czy modele embeddingowe oparte na enkoderach modeli ASR są w stanie dokładnie enkodować semantykę danych wejściowych?}
Powiązana z poprzednim pytaniem, hipoteza ta została potwierdzona w ramach eksperymentów opisanych w Rozdziale 4. Obie proponowane w tej pracy architektury wykazały się stabilnością i wysokimi wynikami maksymalizowanych metryk w ramach treningu. Pomijając wymaganie obecności modelu w ramach całościowego systemu wyszukiwania informacji i rozpatrując zarówno \textit{Attention Pool Model} i \textit{Embedding Alignment Model} w izolacji, stanowią one obiecujące metody tworzenia multimodalnych modeli embeddingowych. Architektura \textit{APM}, celem której treningu jest maksymalizacja podobieństwa między wektorami generowanymi przez gałąź przetwarzania sygnału audio i analogicznymi wektorami dla transkrypcji tych nagrań generowanymi przez pretrenowany model embeddingowy skutecznie dokonuje transformacji między dwoma przestrzeniami embeddingowymi, oferując szybsze działanie modelu niż w przypadku podejścia kaskadowego (Rozdział 5). Model \textit{EAM} przypomina za to bardziej pełnoprawny model embeddingowy, trening którego wymaga potencjalnie ogromnej ilości danych, lecz pozwala na wykorzystanie istniejących modeli pełniących funkcję enkoderów w ramach całościowego systemu. Dodatkową zaletą jest potencjał wykorzystania treningu wstępnego nienadzorowanego, potencjalnie pozwalając na ograniczenie wielkości wymaganych danych zawierających etykiety (pary nagranie mowy -- transkrypcja), które w ramach tej pracy nie poprawiło jakości modelu i trening metodą Data2Vec był niestabilny -- nie świadczy to jednak jednoznacznie o braku przydatności tej metody, lecz w celu wykorzystania tego procesu, należałoby powtórzyć eksperyment z większą ilością danych (Rozdział 4).
\subsection{Jak jakość wyszukiwania opartego na reprezentacjach ukrytych wypada w porównaniu z podejściem kaskadowym (ASR $\rightarrow$ embedding tekstowy)?}
Jakość wyszukiwania jest pojęciem zależnym od oczekiwań użytkownika docelowego całościowego systemu. W Rozdziale 5 przedstawiono szereg porównań i metryk, pozwalających na kompleksowe podejście do wyboru rozwiązania w zależności od wymagań użytkowników systemu. Zaproponowane w tej pracy modele \textit{Attention Pool Model} i \textit{Embedding Alignment Model} cechują się większą prędkością przetwarzania niż klasyczne podejście kaskadowe oparte na transkrypcji, którego ogromną zaletą jest większa wyjaśnialność poprzez produkcję transkrypcji zapytania audio, co potwierdza hipotezę zaproponowaną we wprowadzeniu pracy. W ramach metryk oceniających dokładność wyszukiwania informacji, \textit{APM} i \textit{EAM} wypadają o wiele gorzej niż ewaluowane kaskadowe potoki przetwarzania. Model SONAR, będący multimodalnym modelem embeddingowym, radził sobie lepiej niż proponowane w tej pracy architektury, jednakże nadal nie był w stanie prześcignąć podejścia kaskadowego wykorzystującego model embeddingowy Jina, nawet dla wysokich wartości błędu WER dla transkrypcji (Rozdział 5, Rysunek \ref{fig:wynikiWer}). Zjawisko to pokazuje, że dobór odpowiedniego modelu odpowiedzialnego za semantyczne osadzenie tekstu jest o wiele ważniejsze niż minimalizacja WER w wykorzystanym modelu ASR generującego transkrypcje zapytań w ramach podejścia kaskadowego.
\subsection{Jaki wpływ na jakość obu podejść ma język zapytania (polski a angielski)?}
Poszukiwania odpowiedzi na to pytanie zostały opisane głównie w Rozdziale 5, jak i częściowo w Rozdziale 3. W ramach eksperymentów ewaluujących zachowanie przestrzeni latentnej modeli ASR, zbadane zostało zjawisko powstania osobnych grup reprezentujących różne języki, w ramach którego pokazano pewną inwariancję enkodowanych reprezentacji w świetle języka przetwarzanej wypowiedzi (z wyjątkiem modelu Whisper, gdzie zaobserwowano wyraźną separację między językiem polskim i angielskim). Obserwacja ta została potwierdzona w ramach eksperymentu wielojęzykowego wyszukiwania informacji dla benchamarku PubmedQA w Rozdziale 5 (Rysunek \ref{fig:wynikiJezykAudio}), gdzie model \textit{APM} uczony na danych w języku polskim, wyraźnie gorzej poradził sobie w ramach wyszukiwania informacji dla zapytania w języku angielskim. Z kolei model \textit{EAM} osiągnął lepsze wyniki dla zapytania w języku angielskim mimo nauki na danych w języku polskim. Wyraźne różnice obecne w ramach modeli wykorzystujących translację przestrzeni latentnych nie zostały zaobserwowane przy podejściu kaskadowym, dla którego jakość wyszukiwania utrzymywała się na podobnie wysokim poziomie (zarówno błąd transkrypcji WER, jak i jakość wyszukiwania samego modelu embeddingowego). Podejście kaskadowe cechuje się lepszymi metrykami i stabilnością dla zadania wyszukiwania wielojęzykowego niż podejście latentne.

\section{Dyskusja}
Podczas formułowania zakresu tej pracy wyłoniły się dwa główne nurty, które zostały w niej poruszone. Pierwszym z nich jest sama ewaluacja różnych podejść wyszukiwania informacji za pomocą zapytania głosowego, opisana w Rozdziale 5. Proponowany system pozwala na kompleksowe badania nad różnymi rozwiązaniami wykorzystującymi dodatkowe metody przetwarzania artefaktów generowanych w ramach łańcucha przetwarzania (m.in. augementacja danych, postprocessing embeddingów, korekcja transkrypcji) na różnych zbiorach danych. Nie jest to narzędzie przydatne w ramach eksperymentów przeprowadzonych w ramach tej pracy, lecz może również stanowić pomoc w dalszych badaniach nad zadaniem wyszukiwania informacji za pomocą zapytania głosowego, nie tylko w bazach specjalistycznych. Wracając do samych eksperymentów, ważnym usprawnieniem byłoby stworzenie dedykowanego zbioru danych w języku polskim zawierającym zapytanie audio, jego transkrypcję i listę rankingową najbardziej powiązanych dokumentów. Aktualne badania oparte na tłumaczeniu maszynowym streszczeń artykułów ze zbioru PubmedQA i zbiorze AdmedVoice stanowią niewątpliwie cenne zbiory pozwalające na częściową ewaluację wyszukiwania, lecz posiadanie wyłącznie jednego dokumentu oznaczonego jako mający znaczenie dla zapytania (ang. \textit{relevant}) stanowi duże ograniczenie w kompleksowej ewaluacji całościowego systemu, potencjalnie zaburzając otrzymane wyniki. Dodatkowo, przez koszty związane z procesem tłumaczenia angielskich tekstów na język polski, jedynie 10\% wszystkich artykułów zostało przetłumaczone, zmniejszając rozmiar zbioru danych, na którym oparta była ewaluacja. Uwzględnienie pełnego zbioru przetłumaczonego na język polski może stanowić duże usprawnienie w świetle dokładności metryk opisujących wyniki wyszukiwania. Niezbędnym elementem wymagającym dalszej ewaluacji jest wpływ modelu embeddingowego na jakość wyszukiwania metodą kaskadową. W tej pracy przetestowano 3 różne modele, w ramach których JinaV4 osiągała najlepsze wyniki dla różnych metod ewaluacji. Wartościowym uzupełnieniem tych badań byłaby dodatkowa ewaluacja kolejnych modeli embeddingowych, w szczególności tych dedykowanych do działania w ramach domeny medycznej. \\

Kolejnym aspektem tej pracy były proponowane architektury modeli wykorzystujących translację przestrzeni latentnych enkoderów modeli ASR, \textit{Attention Pool Model} i \textit{Embedding Alignment Model}. W ramach treningu na zbiorze danych AdmedVoice, obie z proponowanych architektur cechowały się stabilnym treningiem i wysokimi wartościami metryk maksymalizowanych w ramach tego procesu na zbiorze testowym. Zgodnie z hipotezą przedstawioną we wprowadzeniu, czas potrzebny do przetworzenia zapytania w formie nagrania mowy i wygenerowania wektora będącego jego reprezentacją był mniejszy niż w podejściu kaskadowym, co stanowi niewątpliwą przewagę tych architektur nad podejściem kaskadowym. Jednakże, w ramach ich całościowej ewaluacji jako część systemu wyszukiwania informacji, nie były one w stanie dorównać przetwarzaniu opartym na generowaniu transkrypcji. Reiterując wnioski powstałe w ramach odpowiedzi na pytania badawcze, najprawdopodobniej winna jest temu zjawisku ilość danych wykorzystanych w ramach ich treningu. Model SONAR, będący multimodalnym modelem embeddingowym, wykazał się większą skutecznością w ramach wyszukiwania informacji, mimo braku dostosowania go do operowania w ramach domeny medycznej. Oznacza to przewagę zdolności generalizacji modelu, uzyskanej na podstawie większej ilości danych treningowych, nad dedykowanym modelem przystosowanym wyłącznie do przetwarzania danych w ramach mowy i tekstów klinicznych. W celu rzetelnego porównania architektur, należałoby wytrenować modele \textit{APM} i \textit{EAM} na o wiele większych zbiorach danych, zawierających również dane spoza domeny medycznej, w celu zwiększenia ich zdolności generalizacyjnej. Praca ta wykazała, że architektury te produkują modele stabilne i zdolne do wyszukiwania informacji, jednakże z tak małą ilością danych obecnych w zbiorze treningowym, nie są one w stanie prześcignąć innych poddanych ewaluacji modeli. Dodatkowo, uczenie modelu \textit{Embedding Alignment Model} z wykorzystaniem treningu wstępnego wymaga dalszej pracy, również na o wiele większym zbiorze danych, w celu zbadania potencjalnych zysków z zastosowania takiego podejścia.
